\section{Related Work}
\label{related_work}

\subsection{Human-in-the-Loop Systems}
\label{subsec:rw-hitl}
Human-in-the-loop systems emphasize continuous user oversight, intervention, and correction. For requirement refinement, the relevant issue is not only whether a human can approve an output, but whether the user can guide revisions during an evolving requirement process.

\subsection{Interactive Requirement Engineering}
\label{subsec:rw-interactive-re}
Interactive requirement engineering studies how stakeholders clarify goals, negotiate constraints, and iteratively improve requirement statements. This work builds on that perspective by treating interaction as a core refinement mechanism rather than an optional interface layer.

\subsection{Requirement Refinement}
\label{subsec:rw-refinement}
Requirement refinement concerns transforming vague, incomplete, or inconsistent requirement drafts into clearer and more actionable specifications. The focus of this paper is limited to requirement-level revision, constraint adjustment, and evolution under user feedback.

\subsection{LLM-Based Requirement Engineering}
\label{subsec:rw-llm-re}
LLM-based requirement engineering explores how generative models can support elicitation, drafting, rewriting, and analysis of requirements. Existing automatic approaches are useful for producing candidate requirements, but they often provide limited support for controlled, transparent, and iterative refinement.

\subsection{Multi-Agent and Human-AI Collaboration}
\label{subsec:rw-multi-agent}
Multi-agent collaboration can distribute responsibilities such as elicitation, analysis, and review. Human-AI collaboration further requires that users remain able to inspect, correct, and redirect the refinement process. This paper combines these two lines of work in a lightweight workflow for interactive requirement optimization.

\subsection{Positioning of This Work}
\label{subsec:rw-positioning}
Unlike work that treats requirements as input to later development tasks, this paper focuses only on requirement refinement itself. The contribution is a research-problem-driven structure for studying how human feedback, agent collaboration, and knowledge assistance improve requirement controllability.
